\documentclass[12pt]{article}
\usepackage[a5paper]{geometry}

% \title{algorithmicx (algpseudocode) example}

% From the package documentation:
% "The package algorithmicx itself doesn’t define any
% algorithmic commands, but gives a set of macros to 
% define such a command set.  You may use only 
% algorithmicx, and define the commands yourself, or you 
% may use one of the predefined command sets."
% 
% Popular predefined command sets include algpseudocode and algpascal. algcompatible should only be used in old documents, and algc is incomplete.

\usepackage{algpseudocode}

\begin{document}

\begin{algorithmic}
\Require $n \geq 0$
\Ensure $y = x^n$
\State $y \Leftarrow 1$
\State $X \Leftarrow x$
\State $N \Leftarrow n$
\While{$N \neq 0$}
\If{$N$ is even}
  \State $X \Leftarrow X \times X$
  \State $N \Leftarrow \frac{N}{2} $  \Comment{This is a comment}
\ElsIf{$N$ is odd}
  \State $y \Leftarrow y \times X$
  \State $N \Leftarrow N - 1$
\EndIf
\EndWhile
\end{algorithmic}



\end{document}